\begin{lemma}[Pointwise slice bounds and the slice gradient for $K_{\gamma,j}$]
Fix $n \geq 1$ and $j \in \{1,\ldots,n\}$, and let
\[
  \Phi_{w,y}(r)
    = \frac{y_j - r w_j}{(1-r^2)^{n/2+1}\sqrt{-\log r}}
      \exp\Bigl(-\frac{|rw-y|^2}{1-r^2}\Bigr),
  \qquad w,y \in \mathbb{R}^n,\ r \in (0,1),
\]
be the $r$-integrand of the normalized Gaussian Riesz kernel of
\noderef{GaussRieszKernel}.  Let $M \geq 3$ and let $y \in \mathbb{R}^n$ satisfy
$|y| \leq M$.  Write $\tau = 1-r$, so that $\tau \in (0,1)$ for $r \in (0,1)$.
Then:
\begin{enumerate}
\item for all $w \in \mathbb{R}^n$ and all $r \in (0,1)$,
\[
  \frac{|rw-y|^2}{1-r^2} = \bigl(|y|^2-|w|^2\bigr) + \frac{|w-ry|^2}{1-r^2} ;
\]
\item for all $w \in \mathbb{R}^n$ with $|w| \leq M$ and all $r \in (0,1)$,
\[
  \bigl|\Phi_{w,y}(r)\bigr|
    \leq e^{M^2}\,\bigl(|w-y| + \tau M\bigr)\,\tau^{-n/2-3/2}
      \exp\Bigl(-\frac{|w-ry|^2}{2\tau}\Bigr) ;
\]
\item there is a map $G$ assigning to each $w \in \mathbb{R}^n$ and each
$r \in (0,1)$ a vector $G(w,r) \in \mathbb{R}^n$ such that
  \begin{enumerate}
  \item for every $w \in \mathbb{R}^n$ the map $r \mapsto G(w,r)$ is continuous
    on $(0,1)$; and
  \item for all $w \in \mathbb{R}^n$ with $|w| \leq M$ and all $r \in (0,1)$,
    the map $v \mapsto \Phi_{v,y}(r)$ is (Fr\'echet) differentiable at $w$ with
    derivative the continuous linear functional
    $h \mapsto \langle G(w,r), h\rangle$, and
    \[
      |G(w,r)| \leq e^{M^2}\,\tau^{-n/2-3/2}
        \Bigl(1 + \frac{2\bigl(|w-y| + \tau M\bigr)^2}{\tau}\Bigr)
        \exp\Bigl(-\frac{|w-ry|^2}{2\tau}\Bigr) .
    \]
  \end{enumerate}
\end{enumerate}

\emph{Why clause (3) produces the gradient rather than bounding a chosen
derivative.}  A bound on a derivative operator selected by convention where a
function is not differentiable is vacuous wherever differentiability fails.
Clause (3)(b) therefore asserts genuine differentiability at $w$ and bounds the
vector that represents the actual derivative there.  Clause (3)(a) records that
the representing vector may be chosen to depend continuously on $r$; this is
what makes $r \mapsto G(w,r)$ measurable on $(0,1)$, which is one of the
hypotheses of the theorem on differentiation under the integral sign that a
consumer applies to $\int_0^1\Phi_{w,y}(r)\,dr$.

\emph{Clause (1) carries no hypothesis on $w$, $y$ or $M$.}  It is an algebraic
identity valid for every $w,y \in \mathbb{R}^n$ and every $r \in (0,1)$; the
hypotheses $M \geq 3$ and $|y| \leq M$ are used only in clauses (2) and (3).
Clause (1) is stated because it is the step that makes the factor
$e^{|w|^2-|y|^2}$ visible, and it is what both of the other clauses are proved
from.

\emph{The bounds are stated for a general $w$, not at a single base point.}  A
consumer of clause (3) needs a bound on $|G(w,r)|$ that is valid simultaneously
at every $w$ of a ball, with the dependence on $w$ exhibited through $|w-y|$,
in order to produce a majorant of $G(\cdot,r)$ that is independent of $w$ on
that ball.  A version pinned at one point would not serve that use.
\end{lemma}
\begin{proof}
Throughout fix $r \in (0,1)$ and write $\tau = 1-r \in (0,1)$,
\[
  c = (1-r^2)^{n/2+1}\sqrt{-\log r},
  \qquad
  A(w,r) = \exp\Bigl(-\frac{|rw-y|^2}{1-r^2}\Bigr) ,
\]
so that $\Phi_{w,y}(r) = (y_j - r w_j)\,A(w,r)/c$.  Note $1-r^2 = \tau(1+r)$, so
$0 < \tau \leq 1-r^2 \leq 2\tau$; and $\log r \leq r-1 = -\tau$, so
$-\log r \geq \tau > 0$.  In particular $\sqrt{-\log r} > 0$ and $c > 0$, and
$A(w,r) > 0$.

\emph{Clause (1).}  Expanding squares,
\[
  |rw-y|^2 = r^2|w|^2 - 2r\langle w,y\rangle + |y|^2,
  \qquad
  |w-ry|^2 = |w|^2 - 2r\langle w,y\rangle + r^2|y|^2 .
\]
Hence
\begin{align*}
  (1-r^2)\bigl(|y|^2-|w|^2\bigr) + |w-ry|^2
    &= |y|^2 - r^2|y|^2 - |w|^2 + r^2|w|^2
       + |w|^2 - 2r\langle w,y\rangle + r^2|y|^2 \\
    &= |y|^2 + r^2|w|^2 - 2r\langle w,y\rangle
     = |rw-y|^2 .
\end{align*}
Dividing by $1-r^2 > 0$ gives clause (1).

\emph{The common estimate $(\ast)$.}  We claim that for $|w| \leq M$ and
$r \in (0,1)$,
\[
  \frac{A(w,r)}{c}
    \leq e^{M^2}\,\tau^{-n/2-3/2}\exp\Bigl(-\frac{|w-ry|^2}{2\tau}\Bigr).
  \tag{$\ast$}
\]
For the numerator, clause (1) gives
$-|rw-y|^2/(1-r^2) = |w|^2 - |y|^2 - |w-ry|^2/(1-r^2)$, so
\[
  A(w,r) = e^{|w|^2-|y|^2}\exp\Bigl(-\frac{|w-ry|^2}{1-r^2}\Bigr) .
\]
Since $|w| \leq M$ and $|y|^2 \geq 0$ we have $|w|^2 - |y|^2 \leq M^2$, so
$e^{|w|^2-|y|^2} \leq e^{M^2}$ by monotonicity of $\exp$.  Since
$0 < 1-r^2 \leq 2\tau$ and $|w-ry|^2 \geq 0$ we have
$|w-ry|^2/(2\tau) \leq |w-ry|^2/(1-r^2)$, hence
$\exp(-|w-ry|^2/(1-r^2)) \leq \exp(-|w-ry|^2/(2\tau))$.  Multiplying the two
(both factors being positive) bounds $A(w,r)$ by
$e^{M^2}\exp(-|w-ry|^2/(2\tau))$.

For the denominator, $0 < \tau \leq 1-r^2$ and $n/2+1 \geq 0$ give
$\tau^{n/2+1} \leq (1-r^2)^{n/2+1}$ by monotonicity of $t \mapsto t^{n/2+1}$ on
$[0,\infty)$; and $0 \leq \tau \leq -\log r$ gives
$\tau^{1/2} = \sqrt\tau \leq \sqrt{-\log r}$.  Both factors on the left being
nonnegative, multiplying the two inequalities gives
$\tau^{n/2+3/2} = \tau^{n/2+1}\tau^{1/2} \leq c$, and $\tau^{n/2+3/2} > 0$, so
$c^{-1} \leq \tau^{-n/2-3/2}$.  Multiplying the numerator and denominator
bounds (all quantities being nonnegative) gives $(\ast)$.

\emph{The two numerator bounds.}  Let $|w| \leq M$.  Since
$y_j - r w_j = (y_j - w_j) + \tau w_j$ and $|y_j - w_j| = |(y-w)_j| \leq |y-w|
= |w-y|$ and $|\tau w_j| = \tau|w_j| \leq \tau|w| \leq \tau M$,
\[
  |y_j - r w_j| \leq |w-y| + \tau M . \tag{6a}
\]
Since $rw - y = (w-y) - \tau w$ and $|\tau w| = \tau|w| \leq \tau M$,
\[
  |rw - y| \leq |w-y| + \tau M . \tag{6b}
\]

\emph{Clause (2).}  Because $c > 0$ and $A(w,r) > 0$,
\[
  \bigl|\Phi_{w,y}(r)\bigr|
    = \bigl|y_j - r w_j\bigr|\cdot\frac{A(w,r)}{c} .
\]
Both factors are nonnegative, and the right-hand sides of $(6a)$ and of
$(\ast)$ are nonnegative, so multiplying $(6a)$ by $(\ast)$ gives clause (2).

\emph{Clause (3): the gradient.}  Fix $r \in (0,1)$.  The map
$v \mapsto y_j - r v_j$ is affine, hence differentiable at every $w$ with
derivative $h \mapsto -r h_j = \langle -r e_j, h\rangle$, where $e_j$ is the
$j$th standard basis vector (so $h_j = \langle e_j,h\rangle$ and $|e_j| = 1$).
The map $v \mapsto rv - y$ is affine with derivative $h \mapsto rh$, so
$v \mapsto |rv-y|^2$ is differentiable at $w$ with derivative
$h \mapsto 2\langle rw-y, rh\rangle = 2r\langle rw-y,h\rangle$; dividing by the
constant $1-r^2$, negating, and composing with $\exp$, the map
$v \mapsto A(v,r)$ is differentiable at $w$ with derivative
\[
  h \mapsto A(w,r)\cdot\Bigl(-\frac{2r}{1-r^2}\Bigr)\langle rw-y,h\rangle .
\]
By the product rule and division by the constant $c$, the map
$v \mapsto \Phi_{v,y}(r)$ is differentiable at every $w \in \mathbb{R}^n$, with
derivative
\[
  h \mapsto \frac{A(w,r)}{c}
    \Bigl[\,-r\langle e_j,h\rangle
      - (y_j - r w_j)\frac{2r}{1-r^2}\langle rw-y,h\rangle\Bigr]
  = \langle G(w,r), h\rangle ,
\]
where we define
\[
  G(w,r) := \frac{A(w,r)}{c}
    \Bigl(-r\,e_j - (y_j - r w_j)\frac{2r}{1-r^2}\,(rw-y)\Bigr) .
\]
This is a continuous linear functional of $h$, being an inner product against a
fixed vector.  So clause (3)(b)'s differentiability assertion holds, with this
$G$, for every $w$ --- in particular for every $w$ with $|w| \leq M$.

\emph{Clause (3)(a): continuity in $r$.}  Fix $w \in \mathbb{R}^n$ and
$r \in (0,1)$.  On $(0,1)$ the functions $r \mapsto 1-r^2$ and
$r \mapsto -\log r$ are continuous and strictly positive (as noted at the
start), so $r \mapsto (1-r^2)^{n/2+1}$ and $r \mapsto \sqrt{-\log r}$ are
continuous and strictly positive, and hence so is $r \mapsto c$.  Also
$r \mapsto |rw-y|^2$, $r \mapsto y_j - r w_j$ and $r \mapsto rw-y$ are
continuous, so $r \mapsto A(w,r)$ is continuous, and $r \mapsto 2r/(1-r^2)$ is
continuous because its denominator does not vanish.  A quotient of continuous
functions with nonvanishing denominator is continuous, and sums, products and
scalar multiples of continuous maps are continuous; therefore
$r \mapsto G(w,r)$ is continuous on $(0,1)$.

\emph{Clause (3)(b): the bound.}  Let $|w| \leq M$.  By the triangle inequality
and homogeneity of the norm,
\[
  |G(w,r)| \leq \frac{A(w,r)}{c}
    \Bigl(r\,|e_j| + |y_j - r w_j|\,\frac{2r}{1-r^2}\,|rw-y|\Bigr) ,
\]
using that $A(w,r)/c > 0$ and $2r/(1-r^2) > 0$.  Now $|e_j| = 1$ and $r \leq 1$,
so the first summand is at most $1$.  For the second, $2r/(1-r^2) \leq 2/\tau$:
indeed both denominators are positive and
$2r\tau \leq 2\tau(1+r) = 2(1-r^2)$ because $r \leq 1+r$.  Combining this with
$(6a)$ and $(6b)$, and using that all the factors involved are nonnegative,
\[
  |y_j - r w_j|\,\frac{2r}{1-r^2}\,|rw-y|
    \leq \bigl(|w-y| + \tau M\bigr)\frac{2}{\tau}\bigl(|w-y| + \tau M\bigr)
    = \frac{2\bigl(|w-y| + \tau M\bigr)^2}{\tau} .
\]
Hence
\[
  |G(w,r)| \leq \frac{A(w,r)}{c}
    \Bigl(1 + \frac{2(|w-y| + \tau M)^2}{\tau}\Bigr) ,
\]
and since the bracket is nonnegative, $(\ast)$ bounds the right-hand side by
\[
  e^{M^2}\,\tau^{-n/2-3/2}
    \Bigl(1 + \frac{2(|w-y| + \tau M)^2}{\tau}\Bigr)
    \exp\Bigl(-\frac{|w-ry|^2}{2\tau}\Bigr) ,
\]
which is clause (3)(b)'s bound.
\end{proof}
